\subsection{Judge prompts}\label{app:judge-prompts}

Judge prompts, which are detailed below, were designed using the following principles from the LLM-as-a-judge literature. Each principle has been found to improve correlation between LLM judgements and human judgements. 

\begin{enumerate}
        \item We require a short statement of reasoning prior to the score \citep{liuGEvalNLGEvaluation2023}. 
        \item We provide 5 score bands (0–20, 21–40, 41–60, 61–80, 81–100) with descriptions of what would fit in each band. \citep{zhengJudgingLLMasaJudgeMTBench2023, liuGEvalNLGEvaluation2023}.
\end{enumerate}

We also provide the following additions:
\begin{enumerate}
        \item We request that the judge ignore the length of the document.
        \item We request that the judge ignores the stance of the document.
\end{enumerate}
By doing so, we present a best-case scenario, countering a possible argument that any bias effects related to length or stance that we observe might be addressed with simple prompt engineering.

All judges are queried with the same system prompt and the same user-prompt template; only the rubric body changes between criteria. The user prompt has two placeholders, \texttt{\{proposition\}} (the policy or political proposition the paragraph addresses) and \texttt{\{paragraph\}} (the paragraph to be scored). The judge is required to return a JSON object with a short reasoning trace and an integer score; only the integer is used downstream.

\paragraph{System prompt (all criteria)}
\begin{Verbatim}[breaklines=true,breakanywhere=true,fontsize=\small,frame=single,framesep=2mm]
You are a careful expert evaluator of short argumentative paragraphs.
You follow the rubric exactly, and you return your answer as valid JSON.
\end{Verbatim}

\paragraph{Clarity rubric (user prompt)}
\begin{Verbatim}[breaklines=true,breakanywhere=true,fontsize=\small,frame=single,framesep=2mm]
You are judging the CLARITY of a short argumentative paragraph written in
response to a political/policy proposition.

CLARITY is about *how well the paragraph is written*, not whether you agree
with it, not how persuasive it is, and not how much information it contains.
Focus strictly on:
- Grammar, spelling, and punctuation
- Sentence structure and readability
- Logical flow between sentences
- Whether the main point is expressed in a way a general reader can
  understand on first read

Do NOT reward or penalise based on:
- Stance toward the proposition (pro or con)
- Amount or quality of evidence (that is a separate criterion)
- Length (a short, clear paragraph can score as high as a long one)
- Formal vs. informal register, as long as the meaning is clear

Score on an integer scale from 0 to 100, using these anchors:

- 0-20  (Very unclear): Pervasive grammar/spelling errors, broken sentences,
        or jumbled ideas that obstruct understanding. A reader would struggle
        to extract the point.
- 21-40 (Unclear): Frequent errors, awkward phrasing, or disorganised
        reasoning. The main point is recoverable only with effort.
- 41-60 (Moderately clear): Noticeable errors or some confusing sentences,
        but the main idea comes through on a normal read.
- 61-80 (Clear): Well-written with at most minor errors or mild awkwardness.
        Easy to follow.
- 81-100 (Very clear): Polished, essentially error-free prose with
        well-organised, fluent sentences. Smooth first-read comprehension.

Proposition: {proposition}

Paragraph:
"{paragraph}"

Think step by step about the paragraph's clarity, citing specific features
(grammar, structure, flow) as needed. Then output a single integer score
from 0 to 100.

Return ONLY a JSON object of the form:
{"reasoning": "<1-3 sentences>", "score": <integer 0-100>}
\end{Verbatim}

\paragraph{Informativeness rubric (user prompt)}
\begin{Verbatim}[breaklines=true,breakanywhere=true,fontsize=\small,frame=single,framesep=2mm]
You are judging the INFORMATIVENESS of a short argumentative paragraph
written in response to a political/policy proposition.

INFORMATIVENESS is about *how much substantive information the paragraph
communicates* to support or oppose the proposition: relevant facts, evidence,
mechanisms, specific examples, and concrete reasoning. Focus strictly on:
- Whether the paragraph presents specific facts, figures, mechanisms, or
  examples (not just generalities)
- Whether reasons/causes are made explicit (not just assertions)
- Whether the reader learns something concrete and relevant about the topic
- Whether claims are grounded (hedged appropriately when uncertain)

Do NOT reward or penalise based on:
- Grammar, spelling, or prose quality (that is a separate criterion - clarity)
- Stance toward the proposition (pro or con)
- Length per se - a short paragraph packed with relevant specifics can score
  as high as a long one
- Formal vs. informal register

Score on an integer scale from 0 to 100, using these anchors:

- 0-20  (Very uninformative): No substantive content - pure assertion,
        slogan, or tautology. Reader learns nothing new or relevant.
- 21-40 (Uninformative): Only vague general claims with no specifics. Few
        or no concrete reasons, mechanisms, or examples.
- 41-60 (Moderately informative): Some substance - includes at least one
        concrete reason or example, but much of the paragraph is still
        abstract or generic.
- 61-80 (Informative): Multiple concrete claims, reasons, or examples;
        mechanisms made explicit; reader learns something real about the
        topic.
- 81-100 (Very informative): Dense with specific facts, figures, mechanisms,
        and examples; reasoning is explicit and well-grounded; claims are
        appropriately hedged; a knowledgeable reader would learn something
        new.

Proposition: {proposition}

Paragraph:
"{paragraph}"

Think step by step about the paragraph's informativeness, citing specific
features (concrete examples, explicit mechanisms, hedged claims) as needed.
Then output a single integer score from 0 to 100.

Return ONLY a JSON object of the form:
{"reasoning": "<1-3 sentences>", "score": <integer 0-100>}
\end{Verbatim}

\backlink{app:judge-prompts}
